\documentclass[12pt,aspectratio=169]{beamer}
\usepackage{amsmath,booktabs,tabularx}

\makeatletter
\def\input@path{{../../../preamble/}{book2026-2027/preamble/}}
\makeatother
\input{graphs/highlights.tex}
% \GTenHighlightsfalse % Uncomment to disable instructional highlights.

\definecolor{LessonInk}{HTML}{132A3E}
\definecolor{LessonTeal}{HTML}{17675E}
\definecolor{LessonTint}{HTML}{EAF2F1}
\setbeamercolor{normal text}{fg=LessonInk,bg=white}
\setbeamercolor{structure}{fg=LessonTeal}
\setbeamercolor{frametitle}{fg=LessonInk}
\setbeamercolor{lesson footer}{fg=LessonTeal,bg=LessonTint}
\setbeamerfont{frametitle}{size=\Large,series=\bfseries}
\setbeamersize{text margin left=9mm,text margin right=9mm}
\setbeamertemplate{navigation symbols}{}
\setbeamertemplate{footline}{%
  \begin{beamercolorbox}[wd=\paperwidth,ht=3ex,dp=1.3ex,leftskip=9mm,rightskip=9mm]{lesson footer}
    \footnotesize\insertsectionhead\hfill\insertframenumber/\inserttotalframenumber
  \end{beamercolorbox}%
}
\setbeamercovered{invisible}
\renewcommand{\arraystretch}{1.22}

\title{Outcomes, Support and Probability Mass Functions}
\subtitle{Grade 10, Term 2, Lesson 1, Criterion C1}
\author{}
\date{}

% Deck-local legend: both frames use this definition; the registry is unchanged.
\newcommand{\LessonHighlightLegend}{%
  \begingroup
  \small
  \setlength{\tabcolsep}{4pt}
  \renewcommand{\arraystretch}{1.0}
  \begin{tabularx}{\linewidth}{@{}>{\raggedright\arraybackslash}p{.255\linewidth}%
    >{\raggedright\arraybackslash}X>{\raggedright\arraybackslash}p{.34\linewidth}@{}}
    \toprule
    \textbf{Highlighted item} & \textbf{What it means} & \textbf{Common symbols} \\
    \midrule
    \GTenHLOutcome{Outcome} & A possible result of the experiment. &
      $\GTenHLOutcome{\omega}$ (one outcome),\newline
      $\GTenHLOutcome{\Omega}$ (sample space). \\
    \GTenHLProbability{Probability} & The chance of an outcome or event, expressed without units. &
      $P(\GTenHLOutcome{\omega})$, $P(\GTenHLValue{X}=\GTenHLValue{x})$;\newline
      e.g.\ result: $\GTenHLProbability{0.4}$. \\
    \GTenHLValue{First variable} & The first variable, its numerical values, and any units. &
      $\GTenHLValue{F}$, $\GTenHLValue{X}$; values $\GTenHLValue{f}$, $\GTenHLValue{x}$. \\
    \GTenHLSecond{Second variable} & The second variable, its numerical values, and any units. &
      $\GTenHLSecond{E}$, $\GTenHLSecond{Y}$; values $\GTenHLSecond{e}$, $\GTenHLSecond{y}$. \\
    \GTenHLThird{Third variable} & The third variable, its numerical values, and any units. &
      $\GTenHLThird{C}$, $\GTenHLThird{Z}$; values $\GTenHLThird{c}$, $\GTenHLThird{z}$. \\
    \GTenHLParameter{Unknown parameter} & A quantity to determine from the information given. &
      $\GTenHLParameter{k}$ (illustrative only): this deck uses no unknown parameter. \\
    \bottomrule
  \end{tabularx}\par
  \endgroup
}

% The three arguments are the first overlays showing the F, E and C columns.
% Reuse the same table when students read the mappings back to form groups.
\newcommand{\ShopMap}[3]{%
  \begin{center}
    \begin{tabular}{@{}lcccc@{}}
      \toprule
      \textbf{Outcome} $\GTenHLOutcome{\omega}$ &
      $P(\GTenHLOutcome{\omega})$ &
      \visible<#1->{$\GTenHLValue{F}(\GTenHLOutcome{\omega})$ (\GTenHLValue{USD})} &
      \visible<#2->{$\GTenHLSecond{E}(\GTenHLOutcome{\omega})$} &
      \visible<#3->{$\GTenHLThird{C}(\GTenHLOutcome{\omega})$} \\
      \midrule
      \GTenHLOutcome{Cash} & \GTenHLProbability{0.1} &
      \visible<#1->{\GTenHLValue{0}} & \visible<#2->{\GTenHLSecond{0}} & \visible<#3->{\GTenHLThird{0}} \\
      \GTenHLOutcome{Debit card} & \GTenHLProbability{0.2} &
      \visible<#1->{\GTenHLValue{1}} & \visible<#2->{\GTenHLSecond{1}} & \visible<#3->{\GTenHLThird{1}} \\
      \GTenHLOutcome{Credit card} & \GTenHLProbability{0.3} &
      \visible<#1->{\GTenHLValue{1}} & \visible<#2->{\GTenHLSecond{1}} & \visible<#3->{\GTenHLThird{1}} \\
      \GTenHLOutcome{Mobile wallet} & \GTenHLProbability{0.4} &
      \visible<#1->{\GTenHLValue{2}} & \visible<#2->{\GTenHLSecond{1}} & \visible<#3->{\GTenHLThird{0}} \\
      \bottomrule
    \end{tabular}
  \end{center}%
}

\newcommand{\ShopFee}{$\GTenHLValue{F}$ is the payment fee in \GTenHLValue{USD}.}
\newcommand{\ShopElectronic}{$\GTenHLSecond{E}$ is $\GTenHLSecond{1}$ for an electronic
  payment and $\GTenHLSecond{0}$ otherwise (unitless).}
\newcommand{\ShopCard}{$\GTenHLThird{C}$ is $\GTenHLThird{1}$ for a card payment
  and $\GTenHLThird{0}$ otherwise (unitless).}

% Reveal the ordered sample space by number of buyers. With all arguments
% equal to 1, the same list is available to consult on the later frames.
\newcommand{\NotebookOutcomes}[4]{%
  \[
    \visible<#1->{\GTenHLOutcome{\Omega}=\{\GTenHLOutcome{\mathrm{NNN}}}%
    \visible<#2->{,\,\GTenHLOutcome{\mathrm{BNN}},\,\GTenHLOutcome{\mathrm{NBN}},\,\GTenHLOutcome{\mathrm{NNB}}}%
    \visible<#3->{,\,\GTenHLOutcome{\mathrm{BBN}},\,\GTenHLOutcome{\mathrm{BNB}},\,\GTenHLOutcome{\mathrm{NBB}}}%
    \visible<#4->{,\,\GTenHLOutcome{\mathrm{BBB}}\}.}
  \]%
}

\newcommand{\NotebookCount}{$\GTenHLValue{X}$ counts the buyers among the three visitors.}
\newcommand{\NotebookRevenue}{$\GTenHLSecond{Y}=5\GTenHLValue{X}$ is total notebook revenue
  in \GTenHLSecond{USD}.}
\newcommand{\NotebookFirst}{$\GTenHLThird{Z}$ is $\GTenHLThird{1}$ if the first visitor buys
  and $\GTenHLThird{0}$ if the first visitor passes (unitless).}
\newcommand{\NotebookDecisions}{In visitor order, \GTenHLOutcome{B} means buy
  (\GTenHLProbability{0.4}) and \GTenHLOutcome{N} means pass (\GTenHLProbability{0.6}).}
\newcommand{\NotebookEventRule}{Independence lets us multiply the decision probabilities
  within an outcome. The matching outcomes are mutually exclusive, so we add their probabilities.}

\newcommand{\NotebookSupports}{%
  \begin{center}
    \begin{tabularx}{\linewidth}{@{}>{\raggedright\arraybackslash}Xll@{}}
      \toprule
      \textbf{Variable} & \textbf{Finite support} & \textbf{Units} \\
      \midrule
      \NotebookCount &
      $\GTenHLValue{\{0,1,2,3\}}$ & \GTenHLValue{buyers} \\
      $\GTenHLSecond{Y}=5\GTenHLValue{X}$ is total notebook revenue. &
      $\GTenHLSecond{\{0,5,10,15\}}$ & \GTenHLSecond{USD} \\
      $\GTenHLThird{Z}$ is $\GTenHLThird{1}$ if the first visitor buys
      and $\GTenHLThird{0}$ if the first visitor passes. &
      $\GTenHLThird{\{0,1\}}$ & Unitless \\
      \bottomrule
    \end{tabularx}
  \end{center}%
}

\begin{document}

\section{Introduction}

\begin{frame}
  \vfill
  {\color{LessonTeal}\bfseries\insertsubtitle\par}
  \medskip
  {\huge\bfseries\inserttitle\par}
  \bigskip
  \textbf{C1.} Characterises discrete random variables in economic and
  financial contexts by identifying their support, probability mass function
  and possible outcomes.
  \vfill
  All data are hypothetical teaching models. Keep exact expressions where
  useful and round final probabilities to four decimal places.
  \vfill
\end{frame}

\begin{frame}[t]{Highlight legend}
  \LessonHighlightLegend
\end{frame}

\begin{frame}[t]{The problem and your solution}
  Given complete, mutually exclusive and exhaustive outcomes, their
  probabilities, and a rule assigning numerical values:

  \medskip
  \textbf{Task 1.} List outcomes and probabilities.\par
  \medskip
  \textbf{Task 2.} Define the variable and units; map outcomes to values;
  give the support and PMF; check validity.\par
  \medskip
  \textbf{Task 3.} Calculate a requested probability and interpret it in context.

  \bigskip
  \begin{center}
    Outcomes $\to$ \shortstack{definition\\and mapping} $\to$
    \shortstack{support\\and PMF} $\to$ checks $\to$
    \shortstack{requested probability\\and interpretation}
  \end{center}
\end{frame}

% Task overview: reveal in task order while reserving all component space.
\begin{frame}[t]{General model: Mapping, PMF and checks}
  \footnotesize
  \setlength{\tabcolsep}{3pt}
  \setlength{\smallskipamount}{1pt}
  \renewcommand{\arraystretch}{1.0}
  \visible<2->{\textbf{Task 2:} $\GTenHLValue{X}$ is [quantity measured];
    units: \GTenHLValue{[units or unitless]}.}

  \smallskip
  \begin{columns}[T,onlytextwidth]
    \begin{column}{.49\textwidth}
      \textbf{Outcome and mapping table}\par
      One row per complete outcome.\par\smallskip
      \begin{tabular}{@{}ccc@{}}
        \toprule
        \multicolumn{2}{c}{\textbf{Task 1}} &
          \visible<2->{\textbf{Task 2}} \\
        \cmidrule(lr){1-2}\cmidrule(l){3-3}
        $\GTenHLOutcome{\omega_i}$ & $q_i=P(\{\omega_i\})$ &
          \visible<2->{$\GTenHLValue{X}(\omega_i)$} \\
        \midrule
        $\GTenHLOutcome{\omega_1}$ & $\GTenHLProbability{q_1}$ &
          \visible<2->{$\GTenHLValue{X}(\omega_1)$} \\
        $\GTenHLOutcome{\omega_2}$ & $\GTenHLProbability{q_2}$ &
          \visible<2->{$\GTenHLValue{X}(\omega_2)$} \\
        $\vdots$ & $\vdots$ & \visible<2->{$\vdots$} \\
        $\GTenHLOutcome{\omega_n}$ & $\GTenHLProbability{q_n}$ &
          \visible<2->{$\GTenHLValue{X}(\omega_n)$} \\
        \bottomrule
      \end{tabular}
    \end{column}
    \begin{column}{.48\textwidth}
      \visible<3->{\textbf{Task 2: PMF (group equal values)}}\par
      \visible<3->{One row per distinct value.}\par\smallskip
      \begin{tabular}{@{}cc@{}}
        \toprule
        \visible<3->{$\GTenHLValue{x_j}$} &
          \visible<3->{$p_X(x_j)=\displaystyle\sum_{i:\,X(\omega_i)=x_j}q_i$} \\
        \midrule
        \visible<3->{$\GTenHLValue{x_1}$} &
          \visible<3->{$\sum_{i:\,X(\omega_i)=x_1}\GTenHLProbability{q_i}$} \\
        \visible<3->{$\vdots$} & \visible<3->{$\vdots$} \\
        \visible<3->{$\GTenHLValue{x_m}$} &
          \visible<3->{$\sum_{i:\,X(\omega_i)=x_m}\GTenHLProbability{q_i}$} \\
        \bottomrule
      \end{tabular}
    \end{column}
  \end{columns}

  \smallskip
  \visible<3->{\textbf{Task 2 --- Support:} $S_X=\{x:p_X(x)>0\}$;\quad
    units: \GTenHLValue{[same as $X$]}.}\par
  \smallskip
  \visible<4->{\textbf{Task 2 --- Checks:} $p_X(x)\geq0$ for every $x$, and
    $\sum_{x\in S_X}p_X(x)=1$.}\par
  \smallskip
  \visible<5->{\textbf{Task 3:} $A$ = [requested values];\quad
    $P(\GTenHLValue{X}\in A)=\sum_{x\in S_X\cap A}p_X(x)$.}\par
  \visible<6->{\textbf{Task 3 --- Interpret:} [The chance that \GTenHLValue{[quantity]}
    is in [requested set] is \GTenHLProbability{[result]}.]}
\end{frame}

\section{School shop}

\begin{frame}[t]{Payment outcomes and their assigned values}
  A school shop receives one payment. The payment method is the
  \emph{complete outcome}. A \emph{random variable} assigns a number
  to each outcome.

  \ShopMap{2}{3}{4}

  \visible<2->{\ShopFee}\par
  \visible<3->{\ShopElectronic}\par
  \visible<4->{\ShopCard}
\end{frame}

\begin{frame}[t]{Supports of the fee and indicators}
  The \emph{support} is the set of values with positive probability.

  \ShopMap{1}{1}{1}

  \ShopFee\par
  \ShopElectronic\par
  \ShopCard\par
  \smallskip
  \textbf{Predict.} Which distinct values appear in each variable's column?

  \smallskip
  \visible<2->{\makebox[\linewidth]{%
    $\GTenHLValue{F}\colon\GTenHLValue{\{0,1,2\}}$ (\GTenHLValue{USD})\quad
    $\GTenHLSecond{E}\colon\GTenHLSecond{\{0,1\}}$ (unitless)\quad
    $\GTenHLThird{C}\colon\GTenHLThird{\{0,1\}}$ (unitless)}}
\end{frame}

\begin{frame}[t]{Probability of a 1 USD fee}
  A \emph{probability mass function} (PMF) gives each value's probability.

  \ShopMap{1}{1}{1}

  \ShopFee\par
  \textbf{Predict.} Which outcomes give $\GTenHLValue{F}=\GTenHLValue{1}$?

  \visible<2->{\GTenHLOutcome{Debit card} and \GTenHLOutcome{Credit card}
    are mutually exclusive, so add their probabilities.}
  \[
    \visible<3->{P(\GTenHLValue{F}=\GTenHLValue{1})
      =\GTenHLProbability{0.2}+\GTenHLProbability{0.3}=\GTenHLProbability{0.5}.}
  \]
\end{frame}

\begin{frame}[t]{Probability of an electronic payment}
  \ShopMap{1}{1}{1}

  \ShopElectronic\par
  \medskip
  \textbf{Predict.} Which outcomes give $\GTenHLSecond{E}=\GTenHLSecond{1}$?

  \visible<2->{\GTenHLOutcome{Debit card}, \GTenHLOutcome{Credit card}
    and \GTenHLOutcome{Mobile wallet}.}
  \[
    \visible<3->{P(\GTenHLSecond{E}=\GTenHLSecond{1})
      =\GTenHLProbability{0.2}+\GTenHLProbability{0.3}
      +\GTenHLProbability{0.4}=\GTenHLProbability{0.9}.}
  \]
\end{frame}

\begin{frame}[t]{Probability of a card payment}
  \ShopMap{1}{1}{1}

  \ShopCard\par
  \medskip
  \textbf{Predict.} Which outcomes give $\GTenHLThird{C}=\GTenHLThird{1}$?

  \visible<2->{\GTenHLOutcome{Debit card} and \GTenHLOutcome{Credit card}.}
  \[
    \visible<3->{P(\GTenHLThird{C}=\GTenHLThird{1})
      =\GTenHLProbability{0.2}+\GTenHLProbability{0.3}=\GTenHLProbability{0.5}.}
  \]
\end{frame}

\begin{frame}[t]{The fee PMF}
  \ShopMap{1}{1}{1}

  \begin{center}
    \begin{tabular}{@{}lccc@{}}
      \toprule
      Fee value $\GTenHLValue{f}$ (\GTenHLValue{USD}) &
      \GTenHLValue{0} & \GTenHLValue{1} & \GTenHLValue{2} \\
      \midrule
      $P(\GTenHLValue{F}=\GTenHLValue{f})$ &
      \GTenHLProbability{0.1} & \GTenHLProbability{0.5} & \GTenHLProbability{0.4} \\
      \bottomrule
    \end{tabular}
  \end{center}

  \emph{Worksheet reading.} See the worked grouping problem for the
  full fee calculation.
\end{frame}

\begin{frame}[t]{The electronic-payment and card-payment PMFs}
  \ShopElectronic\par
  \ShopCard

  \begin{center}
    \begin{tabular}{@{}lcc@{}}
      \toprule
      Electronic-payment value $\GTenHLSecond{e}$ &
      \GTenHLSecond{0} & \GTenHLSecond{1} \\
      \midrule
      $P(\GTenHLSecond{E}=\GTenHLSecond{e})$ &
      \GTenHLProbability{0.1} & \GTenHLProbability{0.9} \\
      \midrule
      Card-payment value $\GTenHLThird{c}$ &
      \GTenHLThird{0} & \GTenHLThird{1} \\
      \midrule
      $P(\GTenHLThird{C}=\GTenHLThird{c})$ &
      \GTenHLProbability{0.5} & \GTenHLProbability{0.5} \\
      \bottomrule
    \end{tabular}
  \end{center}

  $\GTenHLSecond{E}$ and $\GTenHLThird{C}$ share a support but have
  different masses. All probabilities are unitless.
\end{frame}

\section{Notebook sale}

\begin{frame}[t]{Three independent visitor decisions}
  A school enterprise offers a 5 USD notebook to each of three visitors.

  \begin{center}
    \begin{tabular}{@{}lcc@{}}
      \toprule
      \textbf{Decision} & \textbf{Symbol} & \textbf{Probability} \\
      \midrule
      Buy  & \GTenHLOutcome{B} & \GTenHLProbability{0.4} \\
      Pass & \GTenHLOutcome{N} & \GTenHLProbability{0.6} \\
      \bottomrule
    \end{tabular}
  \end{center}

  Assume the visitors decide \emph{independently}. One visitor's decision
  leaves the other visitors' probabilities the same.
\end{frame}

\begin{frame}[t]{Eight ordered outcomes}
  Choose \GTenHLOutcome{B} or \GTenHLOutcome{N} for each visitor.
  Read the three positions in visitor order.

  \medskip
  \textbf{Predict.} List the complete outcomes with zero, one, two and
  three buyers.

  \NotebookOutcomes{2}{3}{4}{5}
\end{frame}

\begin{frame}[t]{Probability that only the first visitor buys}
  \NotebookOutcomes{1}{1}{1}{1}

  In \GTenHLOutcome{BNN}, only the first visitor buys. Independence lets us
  multiply the three decision probabilities in visitor order.

  \medskip
  \textbf{Predict.} Which factor belongs to each position in \GTenHLOutcome{BNN}?
  \[
    P(\GTenHLOutcome{\mathrm{BNN}})
    \visible<2->{=\GTenHLProbability{0.4}}\,
    \visible<3->{\times\GTenHLProbability{0.6}}\,
    \visible<4->{\times\GTenHLProbability{0.6}}\,
    \visible<5->{=\GTenHLProbability{0.144}.}
  \]
\end{frame}

\begin{frame}[t]{Outcomes with the same probability product}
  \NotebookOutcomes{1}{1}{1}{1}

  \textbf{Predict.} Which outcomes have the same factors, possibly in a
  different order?

  \begin{center}
    \begin{tabular}{@{}lc@{}}
      \toprule
      \textbf{Complete outcomes} & \textbf{Probability of each outcome} \\
      \midrule
      \visible<2->{\GTenHLOutcome{NNN}} &
      \visible<3->{$\GTenHLProbability{0.6}^{3}=\GTenHLProbability{0.216}$} \\
      \visible<2->{\GTenHLOutcome{BNN}, \GTenHLOutcome{NBN}, \GTenHLOutcome{NNB}} &
      \visible<3->{$\GTenHLProbability{0.4}\times\GTenHLProbability{0.6}
        \times\GTenHLProbability{0.6}=\GTenHLProbability{0.144}$} \\
      \visible<2->{\GTenHLOutcome{BBN}, \GTenHLOutcome{BNB}, \GTenHLOutcome{NBB}} &
      \visible<3->{$\GTenHLProbability{0.4}^{2}\times\GTenHLProbability{0.6}
        =\GTenHLProbability{0.096}$} \\
      \visible<2->{\GTenHLOutcome{BBB}} &
      \visible<3->{$\GTenHLProbability{0.4}^{3}=\GTenHLProbability{0.064}$} \\
      \bottomrule
    \end{tabular}
  \end{center}
\end{frame}

\begin{frame}[t]{Buyer count, revenue and first visitor}
  \NotebookOutcomes{1}{1}{1}{1}

  \NotebookSupports
\end{frame}

\begin{frame}[t]{Comparing \GTenHLOutcome{BNN} and \GTenHLOutcome{NBN}}
  In visitor order, \GTenHLOutcome{B} means buy and \GTenHLOutcome{N} means pass.\par
  \NotebookCount\par
  \NotebookRevenue\par
  \NotebookFirst\par
  \medskip
  \textbf{Predict.} Which values agree for these two outcomes?

  \begin{center}
    \begin{tabular}{@{}lccc@{}}
      \toprule
      \textbf{Outcome} & $\GTenHLValue{X}$ (\GTenHLValue{buyers}) &
      $\GTenHLSecond{Y}$ (\GTenHLSecond{USD}) & $\GTenHLThird{Z}$ (unitless) \\
      \midrule
      \GTenHLOutcome{BNN} & \visible<2->{\GTenHLValue{1}} &
      \visible<2->{\GTenHLSecond{5}} & \visible<3->{\GTenHLThird{1}} \\
      \GTenHLOutcome{NBN} & \visible<2->{\GTenHLValue{1}} &
      \visible<2->{\GTenHLSecond{5}} & \visible<3->{\GTenHLThird{0}} \\
      \bottomrule
    \end{tabular}
  \end{center}

  \visible<2->{Both outcomes have one buyer, so the count and revenue agree.}\par
  \medskip
  \visible<3->{The first visitor buys in \GTenHLOutcome{BNN} and passes in
    \GTenHLOutcome{NBN}, so $\GTenHLThird{Z}$ differs.}
\end{frame}

\begin{frame}[t]{Probability of exactly one buyer}
  \NotebookCount\par
  \NotebookDecisions
  \NotebookOutcomes{1}{1}{1}{1}

  \textbf{Predict.} Which outcomes give $\GTenHLValue{X}=\GTenHLValue{1}$?

  \visible<2->{Add the probabilities of \GTenHLOutcome{BNN},
    \GTenHLOutcome{NBN} and \GTenHLOutcome{NNB}. These complete outcomes
    are mutually exclusive.}
  \[
    \begin{aligned}
      \visible<3->{P(\GTenHLValue{X}=\GTenHLValue{1})}
      &\visible<3->{=P(\GTenHLOutcome{\mathrm{BNN}})+P(\GTenHLOutcome{\mathrm{NBN}})
        +P(\GTenHLOutcome{\mathrm{NNB}})}\\
      &\visible<4->{=\GTenHLProbability{0.144}+\GTenHLProbability{0.144}
        +\GTenHLProbability{0.144}=\GTenHLProbability{0.432}.}
    \end{aligned}
  \]

  \visible<4->{\emph{Worksheet reading.} The worked school-sale problem
    completes all three PMFs. The notation reference shows how to write
    and check the sums.}
\end{frame}

\begin{frame}[t]{Probability of 10 USD in revenue}
  \NotebookCount\par
  \NotebookRevenue\par
  \NotebookDecisions
  \NotebookOutcomes{1}{1}{1}{1}

  \textbf{Predict.} Which outcomes give $\GTenHLSecond{Y}=\GTenHLSecond{10}$?

  \visible<2->{At 5 USD per notebook, 10 USD means $10\div5=2$ buyers.
    The matches are \GTenHLOutcome{BBN}, \GTenHLOutcome{BNB} and \GTenHLOutcome{NBB}.}
  \[
    \begin{aligned}
      \visible<3->{P(\GTenHLSecond{Y}=\GTenHLSecond{10})}
      &\visible<3->{=P(\GTenHLOutcome{\mathrm{BBN}})+P(\GTenHLOutcome{\mathrm{BNB}})
        +P(\GTenHLOutcome{\mathrm{NBB}})}\\
      &\visible<4->{=\GTenHLProbability{0.096}+\GTenHLProbability{0.096}
        +\GTenHLProbability{0.096}}
        \visible<5->{=\GTenHLProbability{0.2880}.}
    \end{aligned}
  \]

  \visible<3->{\NotebookEventRule}
\end{frame}

\begin{frame}[t]{Probability for the first visitor}
  \NotebookFirst\par
  \NotebookDecisions
  \NotebookOutcomes{1}{1}{1}{1}

  \textbf{Predict.} Which outcomes give $\GTenHLThird{Z}=\GTenHLThird{1}$?

  \visible<2->{Only the first position must be \GTenHLOutcome{B}. The other
    two visitors may buy or pass. The matches are \GTenHLOutcome{BNN},
    \GTenHLOutcome{BBN}, \GTenHLOutcome{BNB} and \GTenHLOutcome{BBB}.}
  \[
    \begin{aligned}
      \visible<3->{P(\GTenHLThird{Z}=\GTenHLThird{1})}
      &\visible<3->{=P(\GTenHLOutcome{\mathrm{BNN}})+P(\GTenHLOutcome{\mathrm{BBN}})
        +P(\GTenHLOutcome{\mathrm{BNB}})+P(\GTenHLOutcome{\mathrm{BBB}})}\\
      &\visible<4->{=\GTenHLProbability{0.144}+\GTenHLProbability{0.096}
        +\GTenHLProbability{0.096}+\GTenHLProbability{0.064}}
        \visible<5->{=\GTenHLProbability{0.4000}.}
    \end{aligned}
  \]

  \visible<3->{\NotebookEventRule}\par
  \medskip
  \visible<6->{This equals the original \GTenHLProbability{0.4} probability
    that the first visitor buys.}
\end{frame}

\section{Discrete and continuous}

\begin{frame}[t]{Discrete supports and purchase time}
  Braces list the values in a support.
  \NotebookSupports

  \visible<2->{A \textbf{countably infinite} support has infinitely many values
    that can be listed one by one, as in $\{1,2,3,\ldots\}$.}\par

  \medskip
  \visible<3->{The \textbf{exact time to complete a purchase} in seconds is continuous.
    It can take values throughout an interval.}\par

  \medskip
  \visible<4->{\textbf{Purchase time recorded in whole seconds} takes discrete values.}
\end{frame}

\section{How to present your solution}

\begin{frame}[t]{The general problem and its three tasks}
  A situation gives \GTenHLOutcome{complete outcomes}, their
  \GTenHLProbability{probabilities}, and information for assigning numbers.
  The outcomes are mutually exclusive (only one occurs) and exhaustive
  (they cover every possibility).

  \medskip
  \textbf{Task 1.} Record every complete outcome and its given probability.

  \medskip
  \visible<2->{\textbf{Task 2.} Define the random variable and its units;
    map every outcome to its value; list the support; group outcomes with
    the same value to build the PMF. Show calculations and verify the
    probability conditions.}

  \medskip
  \visible<3->{\textbf{Task 3.} Use the PMF to find a requested probability
    and interpret it in context.}

  \medskip
  \visible<3->{Several complete outcomes may produce the same numerical value.}
\end{frame}

\begin{frame}[t]{Task 1: Record outcomes and probabilities}
  Use one row for \emph{each complete outcome}, even when some outcomes
  will later have the same numerical value.

  \begin{center}
    \begin{tabular}{@{}lc@{}}
      \toprule
      \textbf{Complete outcome} & \textbf{Probability} \\
      \midrule
      $\GTenHLOutcome{\omega_1}$ & $\GTenHLProbability{q_1}$ \\
      $\GTenHLOutcome{\omega_2}$ & $\GTenHLProbability{q_2}$ \\
      $\vdots$ & $\vdots$ \\
      $\GTenHLOutcome{\omega_n}$ & $\GTenHLProbability{q_n}$ \\
      \bottomrule
    \end{tabular}
  \end{center}

  Define $q_i=P(\{\omega_i\})$: the probability of the single outcome $\omega_i$.
  \[
    \visible<2->{q_i\geq0\quad\text{for every }i,
      \qquad\sum_{i=1}^{n}q_i=1.}
  \]
\end{frame}

\begin{frame}[t]{Task 2: Define the variable and map every outcome}
  Define $\GTenHLValue{X}$ in words: what numerical quantity does it measure?
  State its units, or say that it is unitless.

  \medskip
  Keep every outcome and its probability. Add the assigned value:
  \begin{center}
    \begin{tabular}{@{}lcc@{}}
      \toprule
      \textbf{Complete outcome} & \textbf{Probability} & $\GTenHLValue{X}(\omega_i)$ \\
      \midrule
      $\GTenHLOutcome{\omega_1}$ & $\GTenHLProbability{q_1}$ &
        \visible<2->{$X(\omega_1)$} \\
      $\GTenHLOutcome{\omega_2}$ & $\GTenHLProbability{q_2}$ &
        \visible<2->{$X(\omega_2)$} \\
      $\vdots$ & $\vdots$ & \visible<2->{$\vdots$} \\
      $\GTenHLOutcome{\omega_n}$ & $\GTenHLProbability{q_n}$ &
        \visible<2->{$X(\omega_n)$} \\
      \bottomrule
    \end{tabular}
  \end{center}

  \visible<2->{Each outcome receives exactly one value.
    Different outcomes can receive the same value.}
\end{frame}

\begin{frame}[t]{Task 2: Group values and write the PMF}
  The outcome table has one row per complete outcome.
  The PMF has one row per \GTenHLValue{distinct value}.

  \begin{center}
    \begin{tabular}{@{}lc@{}}
      \toprule
      \textbf{Distinct value} $x$ & $p_X(x)=P(X=x)$ \\
      \midrule
      $\GTenHLValue{x_1}$ & \visible<2->{$\sum_{i:\,X(\omega_i)=x_1}q_i$} \\
      $\vdots$ & \visible<2->{$\vdots$} \\
      $\GTenHLValue{x_m}$ & \visible<2->{$\sum_{i:\,X(\omega_i)=x_m}q_i$} \\
      \bottomrule
    \end{tabular}
  \end{center}

  \visible<2->{Group: $p_X(x)=\sum_{i:\,X(\omega_i)=x}q_i$.
    Add all probabilities mapped to $x$.}\par
  \medskip
  \visible<3->{Support: $S_X=\{x:p_X(x)>0\}$; state units separately.
    If every $q_i>0$, include all distinct mapped values.}\par
  \medskip
  \visible<4->{Verify: $p_X(x)\geq0$ and $\sum_{x\in S_X}p_X(x)=1$.}
\end{frame}

\begin{frame}[t]{Worked model: Buying a carton of fruit}
  A food stall buys one carton of fruit. Exactly one supplier fills the
  order. The probabilities for the orchard, cooperative, wholesaler, and
  grocery are 0.28, 0.19, 0.31, and 0.22, respectively. Their prices for the
  carton are 14, 18, 14, and 22 USD, in the same order.

  \medskip
  \visible<2->{Define $\GTenHLThird{C}$ as the purchase cost of one carton,
    in \GTenHLThird{USD}. A complete outcome identifies the supplier.}

  \medskip
  \visible<2->{\textbf{Predict.} Which suppliers produce the same cost?
    How many rows will the outcome table and the PMF need?}

  \medskip
  \visible<3->{The \GTenHLOutcome{orchard} and \GTenHLOutcome{wholesaler}
    both charge \GTenHLThird{14 USD}. Keep four outcomes, but group
    them into three distinct costs.}
\end{frame}

\begin{frame}[t]{Fruit carton: Mapping, PMF and checks}
  $\GTenHLThird{C}$ is the purchase cost of one carton, in \GTenHLThird{USD}.

  \medskip
  \begin{columns}[T,onlytextwidth]
    \begin{column}{.49\textwidth}
      \textbf{Outcome and mapping table}\par\smallskip
      \begin{tabular}{@{}lcc@{}}
        \toprule
        \textbf{Supplier} & \textbf{Probability} & $\GTenHLThird{C}$ \\
        \midrule
        \GTenHLOutcome{Orchard} & \GTenHLProbability{0.28} & \visible<2->{\GTenHLThird{14}} \\
        \GTenHLOutcome{Cooperative} & \GTenHLProbability{0.19} & \visible<2->{\GTenHLThird{18}} \\
        \GTenHLOutcome{Wholesaler} & \GTenHLProbability{0.31} & \visible<2->{\GTenHLThird{14}} \\
        \GTenHLOutcome{Grocery} & \GTenHLProbability{0.22} & \visible<2->{\GTenHLThird{22}} \\
        \bottomrule
      \end{tabular}
    \end{column}
    \begin{column}{.48\textwidth}
      \textbf{PMF: group equal costs}\par\smallskip
      \begin{tabular}{@{}lc@{}}
        \toprule
        $\GTenHLThird{c}$ (\GTenHLThird{USD}) & $p_C(c)=P(C=c)$ \\
        \midrule
        \visible<3->{\GTenHLThird{14}} & \visible<3->{$0.28+0.31=\GTenHLProbability{0.59}$} \\
        \visible<3->{\GTenHLThird{18}} & \visible<3->{\GTenHLProbability{0.19}} \\
        \visible<3->{\GTenHLThird{22}} & \visible<3->{\GTenHLProbability{0.22}} \\
        \bottomrule
      \end{tabular}
    \end{column}
  \end{columns}

  \smallskip
  \visible<3->{Support: $S_C=\GTenHLThird{\{14,18,22\}}$;
    values are measured in \GTenHLThird{USD}.}

  \visible<4->{All masses are nonnegative, and}
  \[
    \visible<4->{\GTenHLProbability{0.59}+\GTenHLProbability{0.19}
      +\GTenHLProbability{0.22}=1.}
  \]
\end{frame}

\begin{frame}[t]{Task 3: Find and interpret a specific probability}
  \begin{enumerate}
    \item Translate the question into an event involving the variable.
    \item Read its mass from the PMF, or add the masses for all matching values.
    \item Write the probability calculation.
    \item Interpret it in context; use a percentage when helpful.
  \end{enumerate}
  For an exact value: $P(X=a)=p_X(a)$.
  \[
    \visible<2->{P(\GTenHLThird{C}=\GTenHLThird{14})
      =p_C(\GTenHLThird{14})=\GTenHLProbability{0.59}.}
  \]
  \visible<3->{The stall pays exactly \GTenHLThird{14 USD} for the carton
    in \GTenHLProbability{59\%} of these orders.}

  \medskip
  \visible<4->{\textbf{Submit:} outcome table; variable definition and units;
    mapping; support; PMF with grouping calculations; probability checks;
    requested probability and contextual interpretation.}
\end{frame}

\begin{frame}[t]{Highlight legend}
  \LessonHighlightLegend
\end{frame}

\begin{frame}[t]{General model: Mapping, PMF and checks}
  $\GTenHLValue{X}$ is [the numerical quantity measured];
  units: \GTenHLValue{[state units or unitless]}.

  \medskip
  \begin{columns}[T,onlytextwidth]
    \begin{column}{.49\textwidth}
      \textbf{Outcome and mapping table}\par\smallskip
      One row per complete outcome.\par\smallskip
      \begin{tabular}{@{}ccc@{}}
        \toprule
        $\GTenHLOutcome{\omega_i}$ & $q_i=P(\{\omega_i\})$ & $\GTenHLValue{X}(\omega_i)$ \\
        \midrule
        $\GTenHLOutcome{\omega_1}$ & $\GTenHLProbability{q_1}$ &
          \visible<2->{$\GTenHLValue{X}(\omega_1)$} \\
        $\GTenHLOutcome{\omega_2}$ & $\GTenHLProbability{q_2}$ &
          \visible<2->{$\GTenHLValue{X}(\omega_2)$} \\
        $\vdots$ & $\vdots$ & \visible<2->{$\vdots$} \\
        $\GTenHLOutcome{\omega_n}$ & $\GTenHLProbability{q_n}$ &
          \visible<2->{$\GTenHLValue{X}(\omega_n)$} \\
        \bottomrule
      \end{tabular}
    \end{column}
    \begin{column}{.48\textwidth}
      \textbf{PMF: group equal values}\par\smallskip
      One row per distinct value.\par\smallskip
      \begin{tabular}{@{}cc@{}}
        \toprule
        $\GTenHLValue{x_j}$ & $p_X(x_j)=\displaystyle\sum_{i:\,X(\omega_i)=x_j}q_i$ \\
        \midrule
        \visible<3->{$\GTenHLValue{x_1}$} &
          \visible<3->{$\sum_{i:\,X(\omega_i)=x_1}\GTenHLProbability{q_i}$} \\
        \visible<3->{$\vdots$} & \visible<3->{$\vdots$} \\
        \visible<3->{$\GTenHLValue{x_m}$} &
          \visible<3->{$\sum_{i:\,X(\omega_i)=x_m}\GTenHLProbability{q_i}$} \\
        \bottomrule
      \end{tabular}
    \end{column}
  \end{columns}

  \smallskip
  \visible<3->{Support: $S_X=\{x:p_X(x)>0\}$;\quad
    units: \GTenHLValue{[same as $X$]}.}

  \medskip
  \visible<4->{Checks: $p_X(x)\geq0$ for every $x$, and
    $\sum_{x\in S_X}p_X(x)=1$.}
\end{frame}

\end{document}
